% NOTE (23 Nov 2026): Rewritten as the method actually followed. Drafted with AI assistance at the owner's request; reword in your own words before submission.
% NOTE (23 Nov 2026): Assumes the supervisor accepts the reframing in RQ-EVIDENCE.md (RQ1 as a development-level comparison, RQ3 as future work, RQ4 as base versus fine-tuned local model). Change Section 3.1 and the Introduction together if the supervisor decides otherwise.
\chapter{Methodology}\label{methodology-proposed-study}

This chapter describes how JadwAI was built and evaluated. The project followed a design-and-evaluate approach: a working prototype was implemented and then measured on synthetic data. Section~\ref{meth-design} states how each research question was answered and how this differs from the original proposal. The later sections describe the interpretation contract, the development corpus, the extraction methods compared, the scoring, the JAI assistant and its fine-tuning, the engineering verification of the scheduler and validator, and the records kept for reproducibility. The measurements themselves are reported in Chapter~\ref{evidence}.

\section{Research design}\label{meth-design}

The study has three parts, kept separate so that one kind of evidence is not read as another.

\begin{enumerate}
\item \textbf{Offline interpretation experiment} (RQ1, RQ2). Two interpretation methods, a rule-based parser and a locally hosted language model, are run on the same fixed set of statements in four language categories. Their outputs are scored against draft expected answers by the same scorer.
\item \textbf{Assistant and fine-tuning evaluation} (RQ4). The base local model and four fine-tuned versions of it answer the same held-out assistant cases. They are compared with automatic checks, latency measurements and a blind side-by-side rating.
\item \textbf{Engineering verification} (supporting question). The solver, the independent validator, the repair diagnosis and the approval and publication gates are tested with deterministic test oracles and measured on synthetic and adjusted university-scale data.
\end{enumerate}

A valid timetable is not evidence of a correct interpretation, and a correct interpretation does not guarantee a good timetable. For this reason interpretation, scheduling and validation are measured separately.

\paragraph{Changes from the proposal.} The proposal planned a confirmatory comparison of several commercial language models on a sealed, independently annotated held-out set, and a participant study comparing manual and AI-assisted entry. Neither was possible within the project. No budget was approved for paid model calls, no second annotator was available, the expected answers were not reviewed by a human, and no ethics approval or participants were obtained. The research questions were therefore answered as follows (Table~\ref{tab:rq-method}).

\begin{table}[h]
\centering
\small
\begin{tabular}{@{}p{0.2\textwidth}p{0.36\textwidth}p{0.36\textwidth}@{}}
\toprule
Question & Method used & Change from the proposal \\
\midrule
RQ1 extraction accuracy & Rule-based parser and local Qwen2.5-7B on 144 development statements, same scorer & Development-level comparison on draft labels, not a confirmatory held-out result \\
RQ2 clarification and unsupported requests & Clarification and unsupported-request scores from the same runs; clarifying behaviour of JAI on held-out cases & Exploratory; small samples \\
RQ3 effort compared with manual entry & Not conducted & Moved to future work (Section~\ref{meth-rq3}) \\
RQ4 model comparison & Base local model against the same model fine-tuned in four rounds, on identical held-out cases & Local base versus fine-tuned model instead of several paid providers; cost per call is zero \\
Engineering question & Tests, seeded faults and measured runs & As planned \\
\bottomrule
\end{tabular}
\caption{How each research question was addressed.}
\label{tab:rq-method}
\end{table}

\paragraph{What the study does not claim.} It does not compare JadwAI's solver with other timetabling systems, it does not claim that the prototype covers what systems such as UniTime cover~\cite{muller2025itc}, and it does not estimate how well any method would perform on unseen statements written by real administrators.

\section{The system under study}\label{meth-system}

The primary user is the timetable administrator. Instructors and students only view their own published timetable. The administrator enters academic data (instructors, rooms, student groups, courses and their lectures, tutorials and laboratories), then enters rules either through forms or in plain English, formal Arabic or Egyptian Arabic. The workflow that the evaluation relies on is:

\begin{enumerate}
\item A typed statement is converted into a \emph{proposal}: zero or more structured rules, plus any questions that must be answered and any parts that cannot be represented.
\item The administrator reviews, corrects or rejects the proposal. Only approved rules are stored and used. A proposal from either the parser or a language model is never applied automatically.
\item The solver (Google OR-Tools CP-SAT~\cite{google_cpsat}) generates a timetable from the approved inputs.
\item An independent validator, written separately from the solver model, recomputes every supported rule on the saved timetable. A timetable is published only if it passes validation against the current approved inputs.
\end{enumerate}

This separation is what makes it possible to use an imperfect interpretation method safely. An interpretation error can lead to a wrong rule being proposed. It cannot lead to an unreviewed rule being used or to a timetable that breaks an approved hard rule being published. Implementation details are given in Chapter~\ref{chap:implementation}.

\section{Supported rules and the interpretation contract}\label{formal-representation-of-supported-constraints}

\paragraph{Supported rules.} The interpretation methods may propose only the following:
\begin{itemize}
\item instructor unavailable or preferred times;
\item room unavailable times;
\item group unavailable times;
\item permitted times for an activity or a group;
\item a weekly teaching-hour limit for an instructor;
\item a change to the global daily hour limits for students or instructors;
\item a lock on an existing assignment.
\end{itemize}
Each rule is either mandatory or preferred. Times are stored in whole or half hours. Structural requirements, such as complete assignment, no room, instructor or group overlaps, eligible instructors and suitable rooms, always apply, whichever way the rules were entered.

\paragraph{Output contract.} Every interpretation is a JSON object with four parts: a \emph{decision}, a list of \emph{proposals}, a list of \emph{clarification codes} and a list of \emph{unsupported codes}. The decision is one of five values:
\begin{itemize}
\item \texttt{ready}: every part of the statement is represented by supported proposals;
\item \texttt{clarify}: a day, time, person, threshold or meaning is missing or ambiguous, and a question must be asked instead of guessing;
\item \texttt{unsupported}: no supported rule can represent the request faithfully;
\item \texttt{conflict}: the request is understood, but its hard conditions cannot all hold;
\item \texttt{mixed}: some parts are supported and others are unresolved or unsupported.
\end{itemize}
Clarification and unsupported codes are short fixed labels, for example \texttt{day\_required}, \texttt{undefined\_early} or \texttt{ambiguous\_instructor}. They identify the question, not its exact wording.

\paragraph{Correctness rule.} A proposal is correct when its rule type, target entity, mandatory or preferred strength, day and times (or numeric value) are all equal to the expected proposal. Rule titles are free text and are not scored, and the order of proposals is ignored. A duplicate proposal counts as an extra prediction. A case is exactly correct only when the decision, the complete set of proposals and both code lists all match.

\section{Development corpus}\label{dataset-construction}

\paragraph{Releases.} Two releases of a synthetic corpus exist.
\begin{itemize}
\item Release \texttt{0.1.0-draft} (25 October 2026) contains 128 statements in 32 intent families, with 96 development and 32 reserved statements.
\item Development bundle \texttt{0.2.0-draft} (1 November 2026) keeps the 96 development statements unchanged and adds 48 statements in 12 new families for a second fictional institution, a design school. The reserved statements were not read, moved or relabelled.
\end{itemize}
All interpretation results in this thesis use the 0.2.0 development bundle.

\paragraph{Content.} The bundle has 144 statements in 36 families. Each family is one intent written in four language categories: English, formal Arabic, Egyptian Arabic and mixed Arabic/English, giving 36 statements per category. The first context (96 statements) is a fictional department with a Sunday--Thursday week. The second (48 statements) is a design school with a Saturday--Thursday week and a 12-hour regular teaching load with overtime allowed. Each context lists instructors with aliases, rooms, groups, activities, the named teaching periods and a fictional saved timetable for lock requests.

The expected decisions are \texttt{ready} for 76 statements, \texttt{clarify} for 44, \texttt{unsupported} for 16 and \texttt{conflict} for 8. The statements cover:
\begin{itemize}
\item clear and compound requests;
\item negation, including Egyptian Arabic forms;
\item missing days or times;
\item ambiguous and unknown names, for example two instructors called Nour, and a person absent from the context;
\item undefined terms such as ``early'';
\item hard contradictions;
\item requests outside the supported rules;
\item minute-precise times;
\item conversion from sessions to teaching hours.
\end{itemize}

\paragraph{Labels.} The statements and expected answers were written with the help of an AI assistant, following annotation guide~v2. The guide defines each decision, the supported proposal boundary and how ambiguity is to be resolved. Every item records its origin and has the status \texttt{draft\_unreviewed}. No human annotation, second annotator or adjudication took place, although the tools for them were built (gold-free review exports, immutable review records, Cohen's $\kappa$ comparison and adjudication forms). The expected answers are therefore one author's interpretation, and results against them are development figures.

\paragraph{Reserved material is not a blind test.} Reserved statements are excluded from default tooling, and access requires a logged reason. However, their author has seen the labels, so they cannot serve as an independent held-out set. A fresh, independently written and annotated set would be needed for a confirmatory study. Families, not single statements, are the unit for splitting and for confidence intervals, because four translations of one intent are not independent observations.

\section{Interpretation methods compared}\label{baselines}

Both methods receive exactly the same input for each statement: the statement text, the context's entity records (identifier, kind, English name and code) and the calendar. Arabic aliases from the context are not supplied, so Arabic names must be matched to English records by the method itself.

\paragraph{Rule-based parser.} JadwAI's built-in parser uses deterministic patterns for English, formal Arabic and Egyptian Arabic, and matches names across scripts by their consonant sounds. It was measured twice. After the first run, its clarification codes were aligned with guide~v2, three cases where it guessed instead of asking were changed, and missing negations were added (Section~\ref{evidence-baseline}). Because these changes were made after looking at development errors, its reported figures are optimistic for this corpus.

\paragraph{Local language model.} The language-model method uses \texttt{qwen2.5:7b-instruct} (Qwen2.5-7B-Instruct, 4-bit quantised) served by Ollama on the same laptop, with no network access needed. The statement is sent through the same code path as the running application (\texttt{extract\_candidate} with the local provider). The prompt carries the same task instructions as the adapter for an external model. Decoding is constrained by a JSON schema, so the reply always has the contract's shape. Two settings of this schema matter:

\begin{itemize}
\item \textbf{Field order.} The model generates fields in the order the schema lists them. An early run that asked for the decision first produced \texttt{unsupported} for almost every statement. The schema now lists proposals first, then clarifications and unsupported codes, and the decision last, so the decision is written after the model has stated what it found.
\item \textbf{Code format.} Codes must be short lower-case identifiers of 3 to 60 letters, digits and underscores, starting with a letter. This prevents free-text explanations inside the code lists.
\end{itemize}

On a small check of 8 statements, the reordering raised correct decisions from 1 to 5. These 8 statements are part of the development corpus, so the change was also tuned on development data. Decoding used temperature~0, seed~42 and an 8{,}192-token context window. Each statement had one attempt with a 600-second time limit; a failed or invalid reply counts as an error, not as a skipped case.

\paragraph{Second prompt: the evaluation contract.} The first full run showed that the application's prompt no longer matched the evaluation contract: it described newer rule types and a parameter field that the contract rejects, and it did not list the guide's clarification codes (Section~\ref{evidence-llm}). The model was therefore run a second time with a prompt written for the contract. This prompt lists only the contract's proposal kinds and rule types, gives each of the 12 clarification and 4 unsupported codes used in guide~v2 with a one-line meaning, and constrains the code lists to these values. The context also includes the source draft that assignment locks refer to. The model, schema order and decoding settings were unchanged; the statement is sent directly to Ollama rather than through the application. Because the code meanings were written from the codes used in the development labels, this run is tuned toward those labels, and its figures are optimistic in the same way as the parser's.

\paragraph{Shared format adapter.} Before scoring, both methods' outputs pass through the same adapter. It removes empty optional fields and writes whole hours as integers (10 rather than 10.0), which the contract requires. It does not change half-hour values or any meaning.

\paragraph{External model adapter.} An adapter for a commercial model (GPT-5.4 mini) was implemented and tested with recorded responses, but it was never called because no budget was approved. Therefore no commercial model is compared.

\section{Scoring and metrics}\label{metrics}

All outputs are scored by the project's offline scorer, which first validates each output against the contract and the context. An output that names an entity not in the context, or breaks the contract, counts as wrong. Missing outputs fail their case, and every case stays in the denominator. The following measures are reported overall and separately for each language category:

\begin{itemize}
\item \textbf{Exact output} (primary): the proportion of statements whose whole output matches the expected one.
\item \textbf{Decision accuracy}: the proportion with the correct decision.
\item \textbf{Proposal precision, recall and F1}: proposals matched one-to-one with expected proposals.
\item \textbf{Field F1}: agreement on individual fields of matched proposals, as a looser diagnostic.
\item \textbf{Clarification F1} and \textbf{unsupported-request F1}: agreement on the code lists (RQ2).
\item \textbf{Latency}: wall-clock seconds per statement, reported as the median.
\end{itemize}

Confidence intervals are 95\% percentile bootstrap intervals with 1{,}000 replicates, resampling whole families (seed~42), so that the four translations of an intent stay together. With 36 families the intervals are wide, and the per-language figures, which rest on 36 statements each, are exploratory. No significance test between the two methods is used as a main result, because the labels are draft and the parser was adjusted on the same data.

\section{The JAI assistant and its fine-tuning}\label{meth-jai}

RQ4 was reframed as a comparison between a locally hosted base model and the same model after fine-tuning. The task is JAI's conversational role: explaining JadwAI, answering from the timetable facts supplied in the prompt, asking one focused question when a request is vague, declining off-topic requests, preparing at most one action, and replying in the user's language.

\paragraph{Training data.} Examples are fictional conversations in English, formal Arabic and Egyptian Arabic, grouped into families (one intent with several wordings). Whole families are assigned either to training or to the held-out set, so that no wording of a held-out family is seen in training. The prompts are built by the same code that builds the live chat prompt, including the matching entries of a verified guide (Section~\ref{evidence-jai}). From round~3 on, the dataset builder refuses any example whose answer states a fact about JadwAI that is not in that example's guide entries. The examples were drafted with an AI assistant and approved in bulk by the project owner after reading a change list, not item by item.

\paragraph{Training.} Each round trained a QLoRA adapter (rank~16) on Qwen2.5-7B-Instruct with the Unsloth library on a free T4 GPU (Google Colab or Kaggle), with the loss computed only on JAI's final answer. Each adapter was merged into Ollama's 4-bit weights. To stop the small weight changes being lost when the weights were quantised again, each tensor was re-quantised to a different format (4-bit blocks to 6-bit, 6-bit to 5-bit). A separate script compared the merged weights with the base and the adapter, to check how much of the trained change was present in every trained matrix. Three rounds were trained (Table~\ref{tab:jai-rounds}): v1 on the original 210 examples, v2 on the same answers with the guide added to the prompts and an extra epoch, and v3 on 243 examples whose answers were rewritten to follow the guide. Round~4 data (317 training and 85 held-out rows) was prepared from the failures found in live use but not trained in time.

\paragraph{Evaluation.} All candidates answered the same held-out cases with the production prompt, at temperature~0 and seed~42, on the same laptop. The cases were:
\begin{itemize}
\item 16 conversation cases and 5 action-choice cases from held-out families;
\item for v1 and v3, three wordings of a new question that should be answered with ``not sure''.
\end{itemize}
Three kinds of measurement were used:
\begin{enumerate}
\item \textbf{Automatic checks}: whether the reply is valid, whether it proposes an action it should not, whether it is in the user's language, whether it is at most 80 words, and whether the action chosen in the action cases is correct. These checks do not judge whether an answer is true.
\item \textbf{Latency}: mean seconds per answer.
\item \textbf{Blind side-by-side rating} of v1 and v3 on the 19 conversation answers. The two answers to each case were shown in random order without model names, and each was scored from 1 to 5 on four criteria: correct, logical, helpful and natural, giving at most 20 per case. The key was kept in a separate file. The rater was the project owner, who had read a short summary of some replies beforehand, so the rating is neither independent nor fully blind.
\end{enumerate}
Finally, the preferred candidate (v3) was used in the running application with questions taken from its weakest cases, to see whether the held-out results held in real multi-turn use.

\section{Scheduling and validation evaluation}\label{scheduling-validation-and-system-behaviour-evaluation}

\paragraph{Instances.} Three kinds of data were used:
\begin{itemize}
\item small hand-built instances, each designed to exercise one rule;
\item a demonstration department workspace of 448 weekly sessions, 93 groups, 123 instructors and 65 rooms, derived from real course structures with invented names;
\item the Winter 2026 university workbook (2{,}711 weekly sessions), entered through the data template.
\end{itemize}
The university data needed documented adjustments before a timetable could exist, and is therefore used to test the mechanism of generating in parts, not to produce a real university timetable.

\paragraph{Settings and measures.} All runs used one laptop (Intel Core i5-1038NG7, 16\,GB memory, CPU only), CP-SAT with eight search workers and seed~42. The measures recorded were:
\begin{itemize}
\item whether every session was placed;
\item the hard violations found by the independent validator (zero is required for publication);
\item room, instructor and group overlaps;
\item the solver outcome (optimal, feasible, infeasible or unknown);
\item wall-clock time.
\end{itemize}
Because parallel search is not deterministic, a different valid timetable can result from the same inputs.

\paragraph{Testing the validator and the gates.} The following checks were used:
\begin{itemize}
\item \textbf{Seeded violations.} Valid timetables were changed to break exactly one rule, and the validator had to report that rule and the sessions it affects.
\item \textbf{Comparison with the solver model.} Given timetables were checked against the solver model by fixing every assignment variable with an equality constraint (not a hint) and solving for feasibility.
\item \textbf{Gate tests.} Automated tests check that an unapproved proposal never reaches the solver input, that an unsupported request creates no rule, that a short time limit on a feasible instance never reports \emph{infeasible}, and that a timetable failing validation cannot be published.
\end{itemize}

\paragraph{Repair diagnosis.} Deliberately impossible inputs were created, such as locked clashes, an instructor unavailable all week, or a daily limit below the teaching load. The diagnosis had to name a minimal set of changes, and a repaired timetable had to pass validation against everything else (Section~\ref{evidence-repair}).

\paragraph{Automated suites.} Backend, frontend and evaluation-harness test suites were run before each reported measurement. Every fault found during evaluation received a regression test.

\section{AI-assisted versus manual entry (RQ3)}\label{meth-rq3}

The proposal planned a within-subject, counterbalanced pilot. Administrators would enter equivalent task lists once with forms and once with typed statements and review. The measures would be completion time, number of corrections, accuracy of the final approved rules against the expected rules, and whether the participant could explain what an approved rule allows. The recording features needed for this exist in the prototype. The study was not run, because no participants or ethics approval were available in the project period. RQ3 is therefore not answered and is listed as future work. No claim about reduced effort is made.

\section{Reproducibility records}\label{reproducibility-records}

Each reported measurement is stored in the repository with the script that produced it, its log and the code commit it ran on (folders under \texttt{thesis/evidence/}, named by date). For the local-model run, every reply is appended to a predictions file as soon as it arrives, together with its latency and token counts. A stopped run therefore continues from the last saved statement. A run note records the commit, the Ollama version, the machine and each restart, and confirms that the interpretation and evaluation code did not change between restarts. The corpus manifest stores hashes of the data, guide, contract and scorer, so the scores can be recomputed from the saved predictions. Fine-tuning runs keep their exported datasets with hashes, the training notebook, the merge check output, every candidate's answers and the rating sheet with its separate key (\texttt{finetuning/}).

\section{Ethics and privacy}\label{ethics-privacy-and-consent}

\begin{itemize}
\item All statements, contexts and training conversations are fictional. The demonstration workspace uses invented names.
\item The language model runs on the same machine as the application. No statement, timetable or personal data was sent to an external AI service during the evaluation.
\item No human participants took part, so no participant data was collected.
\item AI tools were used to draft the corpus, the training examples and parts of the implementation. This is stated wherever such material is evaluated, and the author is responsible for checking the claims and the records.
\end{itemize}

\section{Threats to validity}\label{threats-to-validity}

\paragraph{Construct validity.} The expected answers encode one interpretation and were not reviewed by a second person. Some statements may have more than one reasonable reading. Exact output is a strict measure, so field F1 is reported alongside it. The automatic checks for JAI measure form (language, length, unwanted actions), not truth, and the human rating had a single rater who was not independent.

\paragraph{Internal validity.}
\begin{itemize}
\item Both interpretation methods were adjusted after looking at development errors: the parser's codes and negations, and the model's schema order. Their scores are therefore optimistic.
\item The same author wrote the corpus, the parser, the prompts and the scorer, and also built the solver and the validator. Shared misunderstandings are possible. Seeded faults and the comparison with the solver model reduce this risk but do not remove it.
\item The AI assistant that helped draft the corpus and training data is itself a language model, and its style may favour models of a similar kind.
\end{itemize}

\paragraph{External validity.}
\begin{itemize}
\item The statements are synthetic and may be easier or harder than real administrators' wording.
\item Egyptian Arabic was written by the author and does not cover regional variation.
\item There are two fictional contexts, one machine, one seed and temperature~0.
\item Results concern one 7-billion-parameter model and may not transfer to other models or later versions.
\end{itemize}

\paragraph{Conclusion validity.} Only 36 families and 16--19 assistant cases were available, so the intervals are wide and small differences between methods or rounds should not be over-read.

\section{Summary of limitations}\label{feasibility-risks-and-limitations}

The evaluation is a development-level study: it uses draft labels, synthetic data, a closed set of supported rules and a single locally hosted model family. It shows how the methods behave on this material and whether the surrounding system prevents unreviewed or invalid output from being used. It does not establish accuracy on unseen real statements, a ranking of commercial models, or any reduction in administrative effort.
